\section{Appendix}
\label{sec:appendix}

\subsection{Per-source breakdown}

\paragraph{Extension set ($2{,}728$ cases).}
F1 by source, budget $r=0.50$, all arms of ours constructed without the answer key. The
routing rule is strongest on \arm{pec\_home} and \arm{smartbench\_tdcs} and weakest on
\arm{toolmem\_raw}, where the uncompressed arm is also weakest. The full breakdown, every arm
against every source, is Table~\ref{tab:persource_all} at the end of this appendix; it is set
full-width because a per-source table with twenty-one arms does not fit a column.


\vspace{2pt}

\paragraph{Main set ($112$ cases).}
Same setting on the smaller set.

\begin{center}
{%
\refstepcounter{table}\label{tab:app_persource_main}%
\footnotesize\textsc{Table \thetable}\\[1pt]
Per-source F1 on the main set ($112$ cases, $r=0.50$).\par}
\vspace{3pt}
{\small\setlength{\tabcolsep}{3pt}%
\resizebox{\columnwidth}{!}{%
\begin{tabular}{lrrrr}
\toprule
Source & $n$ & \arm{full} & \textsc{DARK} & \arm{alias} \\
\midrule
\arm{pec\_home}        & 20 & 0.8750 & 0.8250 & 0.8250 \\
\arm{privacylens}      & 20 & 0.6567 & 0.6400 & 0.4900 \\
\arm{simuhome}         & 20 & 0.8500 & 0.7500 & 0.7500 \\
\arm{smartbench\_tdcs} & 20 & 0.7000 & 0.7000 & 0.6500 \\
\arm{toolmem}          & 12 & 0.4901 & 0.4167 & 0.4354 \\
\arm{toolmem\_raw}     & 20 & 0.6500 & 0.4000 & 0.6500 \\
\midrule
overall                & 112 & 0.7189 & 0.6366 & 0.6475 \\
\bottomrule
\end{tabular}}}
\end{center}
\vspace{2pt}

\subsection{Anchor types covered}

The anchor key is resolved differently depending on what the question names. All six forms
occurring in the extension set are handled, and the anchor string is present in the context
in every case ($2{,}728/2{,}728$).

\begin{center}
{%
\refstepcounter{table}\label{tab:app_anchor}%
\footnotesize\textsc{Table \thetable}\\[1pt]
Anchor key forms covered by the routing rule. All six occur in the extension set and the anchor string is present in the context in every case ($2{,}728/2{,}728$).\par}
\vspace{3pt}
{\small\setlength{\tabcolsep}{3pt}%
\resizebox{\columnwidth}{!}{%
\begin{tabular}{llr}
\toprule
\arm{anchor\_kind} & Example question form & $n$ \\
\midrule
none (entity field) & \emph{for account} \texttt{ACCT-\ldots} & 1500 \\
\arm{id}            & \emph{for record id} \texttt{pec\_home\_\ldots} & 316 \\
\arm{name}          & names an entity directly & 300 \\
\arm{session}       & names a session & 300 \\
\arm{entry}         & \emph{for entry 3} & 281 \\
\arm{native}        & native identifier & 31 \\
\bottomrule
\end{tabular}}}
\end{center}
\vspace{2pt}

\subsection{Leakage counting}

A case counts as leaking when any protected value, or a fragment of it, is recoverable from
the concatenation of the rewritten question and the compressed context. Fragments are the
pattern set of the published protocol. The count is per case, not per value, so the reported
rate is the fraction of cases with at least one exposure; the per-value count is retained in
the artifacts.

\subsection{Reproduction notes}

Every table in this paper is computed from per-case records that ship with the artifacts, one
line per (case, arm). The uncompressed arm in each set is used as an identity check: a result
set whose uncompressed arm does not reproduce the reference value is not part of this paper.
Budget arms assert non-overshoot per case; a run that violates it is discarded rather than
truncated. All arms of ours are keyed alias runs; an unkeyed assignment changes both the
alias strings and the resulting score, so results are only comparable within one setting.

\begin{table*}[t]
\caption{Per-source F1 for every arm on the extension set ($r=0.50$). The aggregate in Table~\ref{tab:external} hides the spread: \textsc{DARK} ranges from $0.2631$ on \arm{toolmem\_raw}, where every arm is weak, to $0.6756$ on \arm{pec\_home}. No source reverses the ordering between \textsc{DARK} and the strongest leak-free published baseline.}
\label{tab:persource_all}
\centering\footnotesize
\setlength{\tabcolsep}{4pt}
\begin{tabular}{lrrrrrrr}
\toprule
Method & \arm{pec\_home} & \arm{privacylens} & \arm{simuhome} & \arm{smartbench\_tdcs} & \arm{toolmem} & \arm{toolmem\_raw} & overall \\
\midrule
\arm{full}                     & 0.7816 & 0.5340 & 0.8686 & 0.6449 & 0.5154 & 0.4401 & 0.6081 \\
\arm{ext\_full}                & 0.7793 & 0.5357 & 0.8686 & 0.6449 & 0.5154 & 0.4384 & 0.6078 \\
\textbf{\textsc{DARK}}         & \textbf{0.6756} & \textbf{0.4671} & \textbf{0.6186} & \textbf{0.5714} & \textbf{0.4933} & \textbf{0.2631} & \textbf{0.4937} \\
Quito-X                        & 0.8505 & 0.1793 & 0.6026 & 0.6163 & 0.4097 & 0.3918 & 0.4831 \\
Quito-X \arm{\_native}         & 0.8520 & 0.1859 & 0.5609 & 0.6408 & 0.4193 & 0.3952 & 0.4862 \\
CPC                            & 0.5523 & 0.4414 & 0.4872 & 0.4571 & 0.4445 & 0.4725 & 0.4740 \\
CPC \arm{\_native}             & 0.5453 & 0.4439 & 0.1218 & 0.5816 & 0.4194 & 0.4772 & 0.4523 \\
LongLLMLingua                  & 0.6864 & 0.3366 & 0.4808 & 0.4401 & 0.2276 & 0.3759 & 0.4230 \\
Provence                       & 0.5630 & 0.4529 & 0.1122 & 0.4755 & 0.3020 & 0.4045 & 0.4080 \\
Provence \arm{\_native}        & 0.5328 & 0.4359 & 0.1058 & 0.4633 & 0.3601 & 0.4126 & 0.4047 \\
LexComp \arm{\_preprint}       & 0.5958 & 0.2719 & 0.0673 & 0.4816 & 0.2912 & 0.4115 & 0.3696 \\
\arm{alias\_budget\_link}      & 0.2859 & 0.3381 & 0.3910 & 0.6367 & 0.3121 & 0.1624 & 0.3482 \\
\arm{budget\_cluster\_link}    & 0.2531 & 0.3342 & 0.3269 & 0.6367 & 0.3154 & 0.1624 & 0.3352 \\
\arm{quad4\_del}               & 0.2742 & 0.4366 & 0.0288 & 0.3776 & 0.4970 & 0.2255 & 0.3142 \\
LLMLingua-2                    & 0.4117 & 0.1784 & 0.0641 & 0.2660 & 0.1668 & 0.4009 & 0.2651 \\
Selective Context              & 0.0215 & 0.1190 & 0.0321 & 0.6388 & 0.2447 & 0.3817 & 0.2588 \\
LLMLingua-2 \arm{\_native}     & 0.2782 & 0.1694 & 0.0545 & 0.2789 & 0.1578 & 0.4041 & 0.2433 \\
\arm{mask\_llmlingua2}         & 0.2709 & 0.2000 & 0.0224 & 0.1418 & 0.1827 & 0.1671 & 0.1712 \\
\arm{mask\_llmlingua2\_native} & 0.2681 & 0.1938 & 0.0214 & 0.1279 & 0.1881 & 0.1622 & 0.1663 \\
\arm{alias\_budget}            & 0.1286 & 0.0668 & 0.3910 & 0.2449 & 0.1626 & 0.0250 & 0.1469 \\
\arm{budget\_cluster}          & 0.1122 & 0.0651 & 0.0160 & 0.2265 & 0.1657 & 0.0250 & 0.0981 \\
\bottomrule
\end{tabular}
\end{table*}

